\section{Funtor $G: Trees \to Struct$}
\label{sec:functor-trees-struct}

\subsection{Descrição das Categorias}

\begin{definition}[Categoria $\mathbf{Trees}$ (modelo MATLAB/Octave)]
Um objeto é uma árvore finita representada por uma estrutura com campos:
\begin{itemize}
    \item \texttt{nodes}: vetor de nós;
    \item \texttt{edges}: matriz $2 \times m$ onde cada coluna é uma aresta $(u, v)$;
    \item \texttt{root}: nó raiz.
\end{itemize}
Um morfismo é uma função entre nós que preserva a estrutura de árvore (ou seja, mapeia a raiz para a raiz e arestas para arestas).
\end{definition}

\begin{lstlisting}[style=matlab, caption={Objeto e morfismo em $\mathbf{Trees}$}]
T.nodes = 1:5;
T.edges = [1 2; 1 3; 2 4; 2 5];
T.root = 1;
\end{lstlisting}

\begin{definition}[Categoria $\mathbf{Struct}$ (modelo MATLAB/Octave)]
Um objeto é um conjunto finito $X$ com uma estrutura adicional (por exemplo, uma relação binária ou uma função). Um morfismo é uma função que preserva a estrutura.
\end{definition}

\begin{lstlisting}[style=matlab, caption={Objeto em $\mathbf{Struct}$}]
S.set = 1:5;
S.rel = [1 2; 1 3; 2 4; 2 5]; % Relação de pai-filho
\end{lstlisting}

\subsection{Funtor Covariante $G: Trees \to Struct$ (Codificação como Conjunto Estruturado)}

O funtor $G: \mathbf{Trees} \to \mathbf{Struct}$ envia cada árvore $T$ a um conjunto estruturado $(X, R)$, onde $X$ é o conjunto de nós de $T$ e $R$ é a relação de pai-filho (ou seja, $(u, v) \in R$ se $u$ é pai de $v$ em $T$). Os morfismos em $\mathbf{Trees}$ são mapeados para morfismos em $\mathbf{Struct}$ que preservam a relação $R$.

\begin{lstlisting}[style=matlab, caption={Funtor $G$ aplicado a uma árvore}]
function S = tree_to_struct(T)
    S.set = T.nodes;
    S.rel = T.edges';
end
\end{lstlisting}

\begin{proposition}
O funtor $G$ satisfaz as seguintes propriedades:
\begin{itemize}
    \item $G(\mathrm{id}_T) = \mathrm{id}_{G(T)}$;
    \item $G$ preserva a composição de morfismos.
\end{itemize}
\end{proposition}

\subsection{Transformações Naturais}

Uma transformação natural $\eta: G \Rightarrow H$ entre funores $G, H: \mathbf{Trees} \to \mathbf{Struct}$ seria uma família de morfismos $\eta_T: G(T) \to H(T)$ que comutam com os morfismos em $\mathbf{Trees}$. Por exemplo, $\eta_T$ poderia ser a identidade se $G = H$.